\section{Experimental protocol}
\label{sec:setup}

\paragraph{Data.}
\emph{PlantVillage}~\cite{hughes2015plantvillage} supplies $54{,}305$ colour images of leaves on a plain background, $38$ classes over $14$ crops. We merge the train and test files of its public Hugging Face copy and resize every image to $128{\times}128$. \emph{PlantDoc}~\cite{singh2020plantdoc} was assembled from internet photographs, so its leaves sit in cluttered scenes. We use the cropped-leaf classification release ($2569$ images, $28$ label names). One label has two images and is dropped, which leaves $27$ classes; images are resized on the short side and centre-cropped to $128{\times}128$. No PlantDoc image is used for training or pruning. \emph{EuroSAT-MS}~\cite{helber2019eurosat} has $27{,}000$ Sentinel-2~\cite{drusch2012sentinel2} patches of $64{\times}64$ pixels (a $640$\,m footprint at $10$\,m resolution) with $13$ spectral bands and $10$ land-use classes. Each band is standardised with the mean and standard deviation of the base-class images. To accept $13$ channels the first convolution repeats the ImageNet filters over the bands and scales them by $3/13$; a control model sees only the three visible bands, B4, B3 and B2, through the unmodified filters.

\paragraph{Splits and episodes.}
A seed fixes a random partition of the classes. On PlantVillage it gives $24$ base, $6$ validation and $8$ novel classes; on EuroSAT-MS, $5$ base and $5$ novel. Seeds $0$ to $4$ give five different partitions, and every method sees the same one within a seed. PlantDoc is an evaluation set only: its $27$ classes are used for $5$-way episodes after training on PlantVillage base classes, so the test crosses a domain boundary. Because PlantDoc and PlantVillage share some concepts (tomato early blight, for one), this is a field-image test and not a strict novel-class test. Every table reports $5$-way accuracy over $2000$ episodes with $15$ query images per class and $K\in\{1,5,10\}$ supports. The episodes are drawn once per dataset, seed and $K$ and reused by all methods, which makes the comparisons paired.

\paragraph{Training and measurement.}
Stage S0 trains for $6$ epochs with SGD (Nesterov momentum $0.9$, learning rate $0.02$, cosine decay, batch $128$, label smoothing $0.1$). Stage R is one epoch at $0.01$. Stage M uses two episodes per step and a learning rate picked on the validation classes of the first partition from $\{10^{-4},2.5{\cdot}10^{-4},5{\cdot}10^{-4},10^{-3},2{\cdot}10^{-3}\}$ (it chose \num{tune:best}). Training images get a random crop (scale $0.6$ to $1$) and random flips; RGB images also get brightness and contrast jitter. Mixed precision is on. On EuroSAT-MS the schedule is the same except for $8$ epochs in S0, $2$ in R and $800$ episodes in M. Apart from that learning rate we tuned nothing, and nothing per method. Off-the-shelf baselines (MobileNetV3-Small, MobileNetV2 and EfficientNet-B0 with torchvision ImageNet weights) get the S0 recipe followed by one epoch of R. All runs used one \num{hw-gpu}. Parameters and multiply-accumulates (MACs) are counted exactly on the pruned network; latency is measured on that GPU and on a server CPU (\num{hw-cpu}) and is not a statement about embedded hardware.

\paragraph{Reporting.}
Tables give the mean over seeds of the per-seed mean accuracy, with the sample standard deviation over the \num{n-seeds} seeds as a subscript. This spread contains the effect of the class partition and of training noise, and it is much larger than the sampling error of $2000$ episodes ($\approx$0.2 points), so we quote it whenever a difference is small. For a pair of methods we give the mean paired difference in points with a $95\%$ $t$-interval over the per-partition differences, and the number of partitions in which the sign agrees. Partitions, not episodes, are the replicates: episodes from one trained network are not independent draws of the training run.
